\documentclass[11pt,a4paper]{article}
\usepackage[a4paper,margin=24mm]{geometry}
\usepackage[T1]{fontenc}
\usepackage[utf8]{inputenc}
\usepackage{lmodern}
\usepackage{amsmath,amssymb,amsthm}
\usepackage{graphicx}
\usepackage{booktabs,array,tabularx,multirow}
\usepackage[font=small,labelfont=bf]{caption}
\usepackage{enumitem}
\usepackage{microtype}
\usepackage{xcolor}
\definecolor{linkblue}{rgb}{0,0.15,0.5}
\usepackage[colorlinks=true,linkcolor=linkblue,citecolor=linkblue,urlcolor=linkblue]{hyperref}
\usepackage{url}
\graphicspath{{./}{figs/}{paper02_probabilistic_sufficiency_v16_src/figs/}}
\theoremstyle{plain}
\newtheorem{theorem}{Theorem}[section]
\newtheorem{proposition}[theorem]{Proposition}
\newtheorem{lemma}[theorem]{Lemma}
\newtheorem{corollary}[theorem]{Corollary}
\theoremstyle{definition}
\newtheorem{definition}[theorem]{Definition}
\newtheorem{example}[theorem]{Example}
\theoremstyle{remark}
\newtheorem{remark}[theorem]{Remark}
\newcommand{\R}{\mathbb{R}}
\newcommand{\Qr}{\mathbb{Q}}
\newcommand{\N}{\mathbb{N}}
\newcommand{\E}{\mathbb{E}}
\newcommand{\Prob}{\mathbb{P}}
\newcommand{\calV}{\mathcal{V}}
\newcommand{\calS}{\mathcal{S}}
\newcommand{\calW}{\mathcal{W}}
\newcommand{\calD}{\mathcal{D}}
\newcommand{\Wfb}{\calW^{\mathrm{fb}}}
\newcommand{\Wbl}{\calW^{\mathrm{bl}}}
\newcommand{\supp}{\operatorname{supp}}
\newcommand{\Post}{\operatorname{Post}}
\newcommand{\hold}{\mathrm{hold}}
\newcommand{\ol}{\mathrm{ol}}
\newcommand{\seq}{\mathrm{seq}}
\newcommand{\obs}{\mathrm{obs}}
\newcommand{\Tobs}{T_{\mathrm{obs}}}
\newcommand{\Tl}{T_{\mathrm{learn}}}
\newcommand{\ind}[1]{\mathbf{1}[#1]}
\newcommand{\tfr}[2]{\tfrac{#1}{#2}}
\setlength{\parskip}{2pt}
\emergencystretch=2em
\allowdisplaybreaks

\begin{document}

\title{\textbf{Probabilistic sufficiency under partial observation:\\ belief-state safety values, survivable-set antichains,\\ and the exact price of late information}}
\author{Amin Abaee\\[0.3em]
{\small Independent Researcher}\\[0.2em]
{\small ORCID: \href{https://orcid.org/0000-0002-0019-1842}{0000-0002-0019-1842}\quad
\href{mailto:amin_abaee@ut.ac.ir}{amin\_abaee@ut.ac.ir}}}
\date{October 15, 2026}
\maketitle

\begin{abstract}
\noindent
Set-valued viability theory under incomplete observation delivers binary verdicts: an information state is viable or it is obstructed. This paper develops the quantitative counterpart, the belief-state safety value $V_k(b)$---the maximal probability, over a declared policy class, of keeping the state in the safe set for $k$ steps when the controller holds a distribution $b$ over states---and carries it entirely in exact rational arithmetic. Three structural results connect the two theories. First, on finite models with deterministic kernels and deterministic observation maps the posterior support follows the set-valued post-state exactly, so $V_k(b)=1$ holds precisely when the support lies in the feedback viable-set recursion, and otherwise the deficit obeys $1-V_k(b)\ge\min_{x\in\supp b}b(x)$; the condition is necessary as well as sufficient, in contrast to the sufficient-only output-feedback criteria of the viability literature. Second, the $\alpha$-vectors of the value are exactly the indicators of the maximal jointly survivable subsets of the support---an antichain, Sperner-bounded, the bound being attained---and the value of observation over blind control is at most $1-\max_x b(x)$, a bound attained for every support size. Third, on a delayed hidden-regime class the hold-class, open-loop and fully observed values are obtained in closed form for all initial states, window lengths, horizons and prior weights; the class lattice orders them, and the two gaps---the value of flexibility and the value of observation---are each $\min(w,1-w)$ on disjoint, explicitly located cells. The classical timing obstruction is thereby shown to be a statement about the declared hold class, exactly on the cells where open-loop time variation restores safety. Beyond the deterministic layer, observation becomes a design variable: a pinned probe of a hidden drift parameter at deadline $T$ moves the viability threshold linearly, $z_0\ge\tfrac{21}{10}+\tfrac{T}{10}$, between a flat free-probe benchmark and a probe-free one, and a general hold--probe--learn schedule has an exact threshold that reduces to an additive law under aligned worst losses. Under noisy sensing the deficit is the prior mass weighted by fatal-reading probabilities, a two-state sensor attains deficit $\tfrac1{20}$ against a minimal mass of $\tfrac12$, a single noisy probe yields exactly three value levels, and $n$ repeated readings bring the wrong-majority probability below $(4\varepsilon(1-\varepsilon))^{n/2}$ at a cost of one unit of floor margin per reading. A contamination operator interpolates exactly between the expected and worst-case readings at one stage, but not across stages, and admits a closed form on the delayed class. Every statement is regenerated from the definitions by a deterministic exact-arithmetic script; the support identity, the stabilization count and the rationality of the $\alpha$-vectors are additionally machine-checked in Lean~4.
\end{abstract}

\medskip
\noindent\textbf{Keywords:} partially observable Markov decision processes; stochastic viability; safety probability; belief states; $\alpha$-vectors; value of information; learning deadlines; exact rational verification

\medskip
\noindent\textbf{MSC 2020:} 93E20, 90C40, 93B07, 49N30, 60J20

